% ID: SEL-D0052-OA-E04
% ESTADO: cotejada. Segunda pasada visual y borrador compilado revisados.
% Fuente: Selectividad/sel_2013_matematicas_aplicadas/septiembre/reserva b/EXAMEN 4-ANDALUCÍA.pdf
% SHA-256: 6201aa5a68f0d1302cfd697751f74c4c73f655f0d9b1ae98ae8960207dbd5225
% Página: 1; opción A; ejercicio original 4.

Se conoce que la acidez de una solución es una variable aleatoria que sigue una distribución Normal con desviación típica 0.2. Se ha tomado una muestra aleatoria de cinco soluciones y se han obtenido las siguientes medidas de la acidez:
\[
7.92\quad7.95\quad7.91\quad7.9\quad7.94.
\]
\begin{enumerate}[label=\alph*)]
\item \textbf{(1.25 puntos)} Halle el intervalo de confianza, al 99\%, para la media poblacional.
\item \textbf{(0.5 puntos)} \textquestiondown{}Qué error máximo se ha cometido en el intervalo anterior?
\item \textbf{(0.75 puntos)} Para el mismo nivel de confianza, calcule el tamaño mínimo muestral que permita reducir el error anterior a la mitad.
\end{enumerate}
